% SLICE-UNSAT REPORT DIRECTIONALTELESCOPE1 2026-10-05 1011 PT
\documentclass{article}
\usepackage{amsmath,amssymb}
\begin{document}
\section*{TASK-008: DIRECTIONALTELESCOPE1}
P versus NP remains open.

\section*{Preflight}
Current AGENTS.md, CODEX.md, the pnp-research skill, root main.tex, TASK-008, DIRECTIONALGAP1, and SLICEDELTA0-R1 were reviewed. Primary literature checked: Cavalar--Kumar--Rossman (Algorithmica 2022 / arXiv:2012.03883) and Cavalar--Fabris--Mukhopadhyay--Srinivasan--Yehudayoff (ECCC TR26-220, 28 September 2026). The former uses gate-by-gate monotone approximation with positive and negative test distributions; the latter gives robust-sunflower control under local pseudo-independence hypotheses. Neither supplies the SAT-specific positive-side bound needed below.

\section*{DIRECTIONAL-TELESCOPE}
Let $h_0,h_1,\ldots,h_L:\Omega\to\{0,1\}$. Let
\[
U=\{i:h_{i-1}\le h_i\text{ pointwise}\},\qquad
D=\{i:h_i\le h_{i-1}\text{ pointwise}\},
\]
and assume every transition belongs to $U\cup D$. Then for arbitrary distributions $\mu_-,\mu_+$,
\[
\boxed{\Pr_{\mu_-}[h_0=0,h_L=1]
\le \sum_{i\in U}\Pr_{\mu_-}[h_{i-1}=0,h_i=1],}
\]
\[
\boxed{\Pr_{\mu_+}[h_0=1,h_L=0]
\le \sum_{i\in D}\Pr_{\mu_+}[h_{i-1}=1,h_i=0].}
\]
\textbf{Status: PROVED/EXACT.}

\paragraph{Proof.}
For each fixed input, a binary path starting at 0 and ending at 1 has a first $0\to1$ transition, which cannot occur at a pointwise-decreasing step. Thus the endpoint error event is contained in the union of the upward-transition error events. Apply the union bound. The second inequality is identical with $1\to0$ and decreasing steps.

\section*{Directional closure/trimming accounting}
If $h_0=f$ and $h_L=A$, and every closure/plucking step is pointwise enlarging while every trimming step is pointwise shrinking, define
\[
c_i=\Pr_{\mu_-}[h_{i-1}=0,h_i=1],\qquad
\tau_i=\Pr_{\mu_+}[h_{i-1}=1,h_i=0].
\]
Then
\[
\boxed{\eta_-\le\sum_{i\in U}c_i,\qquad
\eta_+\le\sum_{i\in D}\tau_i.}
\]
Together with DIRECTIONALGAP1 and SLICEDELTA0-R1 at $t=\lfloor3.5q\rfloor$, a sufficient remaining contradiction budget is
\[
\boxed{\sum c_i+\sum \tau_i+\varepsilon_A<1-o(1).}
\]
This is accounting only. It does not prove local error bounds or moment fooling.

\section*{Falsification-first consequence}
The lemma applies unchanged to XOR-SLICE-EASY, so it is not SAT-specific. It exposes a quantitative failure mode: with polynomially many approximation steps, constant per-step error is useless under this union bound. The next target is \textbf{SAT-LOCAL-ERROR1}: prove a summably small total local error, or another explicit total bound leaving room for $\varepsilon_A$, for every relevant small monotone circuit under the declared SAT distributions. \textbf{OPEN}.

Negative-side closure errors may be candidates for robust-sunflower or pseudo-independence control. Positive-side trimming errors require SAT-specific UNSAT minterm/core geometry or another theorem. No such theorem is claimed here.

\section*{Guards}
XOR-SLICE-EASY, LOCALPATTERN-GUARD, the P/poly positive-family barrier, SLICELIFT2, and THRESHOLD-APPROX-GUARD remain active. The withdrawn switching-lemma route remains \textbf{FALSIFIED/RETIRED}. No restricted-model lower bound is transferred to unrestricted SAT runtime.

\section*{Statuses}
DIRECTIONAL-TELESCOPE: \textbf{PROVED/EXACT}. SAT-LOCAL-ERROR1: \textbf{OPEN}. SAT-SUNFLOWER2-DIRECTIONAL: \textbf{OPEN}. MATCHED-MUF2 and growing-order UNSAT moment camouflage: \textbf{OPEN}. P versus NP remains open.

No Python computation was used.

% END OF VERIFIED REPORT DIRECTIONALTELESCOPE1
\end{document}
